\chapter{1986:Dudley R. Herschbach and Yuan T. Lee and John C. Polanyi}

\label{chap:chemistry1986}

The Nobel Prize in Chemistry for the year 1986 was awarded jointly to Dudley R. Herschbach and Yuan T. Lee and John C. Polanyi.

\textbf{Motivation:}

for their contributions concerning the dynamics of chemical elementary processes

\section{Dudley R. Herschbach}

\subsection*{Early Life and Education}

Dudley R. Herschbach was born in 1932 in San Jose, California, USA.

\subsection*{Career and Major Research}

Herschbach studied at Stanford University (B.S., 1954; M.S., 1955) and earned his Ph.D. in chemical physics from Harvard University in 1958. He taught at the University of California, Berkeley, before joining Harvard University in 1963, where he remained for his career.

\subsection*{Nobel Prize-winning Work}

The work for which Dudley R. Herschbach was awarded the Nobel Prize in Chemistry in 1986 was described by the Nobel Committee as: "for their contributions concerning the dynamics of chemical elementary processes".

He developed the crossed molecular beam technique, in which beams of reactant molecules are directed at one another to study individual reactive collisions and the forces that govern them.

\subsection*{Significance and Impact}

Crossed molecular beams provided the first detailed experimental view of how reactant energy and orientation determine the outcome of elementary chemical reactions. The method launched the modern field of reaction dynamics.

\subsection*{Later Career and Honors}

Herschbach continued teaching and research at Harvard University, where he is professor emeritus. He shared the 1986 Nobel Prize in Chemistry and was awarded the National Medal of Science in 1991.

\subsection*{Legacy}

His crossed-beam experiments, extended by Yuan T. Lee, remain a standard approach for studying elementary reaction mechanisms.

\subsection*{Selected Publications}

"Molecular Dynamics of Elementary Chemical Reactions" (Science, 1987)

\clearpage

\section{Yuan T. Lee}

\subsection*{Early Life and Education}

Yuan T. Lee was born in 1936 in Hsinchu, Taiwan.

\subsection*{Career and Major Research}

Lee received his B.S. from National Taiwan University (1959) and his Ph.D. from the University of California, Berkeley (1965), where he studied under Dudley R. Herschbach. He served on the faculty of the University of Chicago before joining the University of California, Berkeley, in 1974.

\subsection*{Nobel Prize-winning Work}

The work for which Yuan T. Lee was awarded the Nobel Prize in Chemistry in 1986 was described by the Nobel Committee as: "for their contributions concerning the dynamics of chemical elementary processes".

He built a universal crossed molecular beam apparatus that could detect virtually every product of a bimolecular reaction, allowing extremely detailed studies of reaction dynamics.

\subsection*{Significance and Impact}

Lee's instrument extended crossed-beam studies to complex reactions, including hydrocarbon combustion and photodissociation. It produced the most complete experimental picture of elementary chemical processes at the time.

\subsection*{Later Career and Honors}

Lee returned to Taiwan in 1994 to serve as president of Academia Sinica until 2006. He shared the 1986 Nobel Prize in Chemistry.

\subsection*{Legacy}

His universal detection method remains a cornerstone of experimental reaction dynamics.

\subsection*{Selected Publications}

"Molecular Beam Studies of Elementary Chemical Processes" (Science, 1987)

\clearpage

\section{John C. Polanyi}

\subsection*{Early Life and Education}

John C. Polanyi was born in 1929 in Berlin, Germany.

\subsection*{Career and Major Research}

Polanyi studied at the University of Manchester (B.Sc., 1949; Ph.D., 1952), then worked at the National Research Council in Ottawa and at Princeton University. He joined the University of Toronto in 1956, where he was appointed University Professor in 1974.

\subsection*{Nobel Prize-winning Work}

The work for which John C. Polanyi was awarded the Nobel Prize in Chemistry in 1986 was described by the Nobel Committee as: "for their contributions concerning the dynamics of chemical elementary processes".

He developed infrared chemiluminescence, a technique that detects the infrared light emitted by newly formed product molecules to reveal how energy is distributed among their vibrations.

\subsection*{Significance and Impact}

Infrared chemiluminescence made it possible to map the energy released in exothermic reactions, providing direct experimental evidence for how energy is partitioned into new chemical bonds.

\subsection*{Later Career and Honors}

Polanyi continued his research at the University of Toronto, and he later studied the dynamics of reactions at surfaces. He shared the 1986 Nobel Prize in Chemistry.

\subsection*{Legacy}

His findings on reaction energy disposal have fundamental importance for chemiluminescence, lasers, and the spectroscopy of reaction products.

\subsection*{Selected Publications}

"Some Concepts in Reaction Dynamics" (Science, 1987)

\clearpage

\section*{Chapter Summary}

The 1986 Nobel Prize in Chemistry was awarded for the study of chemical elementary processes at the molecular level. Herschbach and Lee provided detailed experimental methods using crossed molecular beams, while Polanyi contributed infrared chemiluminescence for measuring reaction energy disposal.
